% Generated draft: requires mathematical and source review.
\documentclass[11pt]{article}
\usepackage[margin=25mm]{geometry}
\usepackage{fontspec}
\setmainfont{Latin Modern Roman}
\usepackage{amsmath,amssymb,mathtools,unicode-math}
\setmathfont{Latin Modern Math}
\usepackage{enumitem,needspace,xurl,hyperref}
\setlength{\parindent}{0pt}
\setlength{\parskip}{5pt}
\setlength{\emergencystretch}{3em}
\newcommand{\N}{\mathbb{N}}
\newcommand{\Z}{\mathbb{Z}}
\newcommand{\Q}{\mathbb{Q}}
\newcommand{\R}{\mathbb{R}}
\newcommand{\C}{\mathbb{C}}
\newcommand{\F}{\mathbb{F}}
\newcommand{\A}{\mathbb{A}}
\newcommand{\PP}{\mathbb{P}}
\newcommand{\E}{\mathbb{E}}
\newcommand{\eps}{\varepsilon}
\newcommand{\dd}{\,\mathrm{d}}
\begin{document}
\section*{UnsolvedMath ID: 30004528}
\textbf{Source problem code:} catalogue entry
\subsection{Definitions and mathematical statement}
The supplied question is: ``For neighborly spheres, does every linear system of parameters have the Lefschetz property?'' Its precise interpretation requires a coefficient field, a class of spheres, and a choice of Lefschetz property, none of which is specified in the supplied statement.

The following gives a conditional algebraic formulation. Let \(k\) be a field and let \(\Delta\) be a finite simplicial complex of dimension \(d-1\) on vertex set \([n]=\{1,\ldots,n\}\). Its Stanley--Reisner ring, with the standard grading, is
\[
 k[\Delta]=k[x_1,\ldots,x_n]/I_\Delta,
 \qquad
 I_\Delta=\left(\prod_{j\in F}x_j : F\subseteq[n],\ F\notin\Delta\right).
\]
A linear system of parameters is a sequence \(\Theta=(\theta_1,\ldots,\theta_d)\) of degree-one elements of \(k[\Delta]\) for which
\[
 A(\Delta,\Theta)=k[\Delta]/(\theta_1,\ldots,\theta_d)
\]
is finite-dimensional over \(k\). Write \(A_i\) for its degree-\(i\) component.

A standard convention calls a \((d-1)\)-dimensional simplicial sphere neighborly if every subset of at most \(\lfloor d/2\rfloor\) vertices is a face. This convention is a candidate interpretation here; the supplied record does not specify its convention or whether spheres are topological, piecewise-linear, or homology spheres.

The weak Lefschetz property for \(A=A(\Delta,\Theta)\) means that there exists \(\ell\in A_1\) such that, for every integer \(i\geq0\),
\[
 \operatorname{rank}_k(\,\cdot\ell:A_i\longrightarrow A_{i+1})
 =\min\{\dim_k A_i,\dim_k A_{i+1}\}.
\]
The strong Lefschetz property means that there exists \(\ell\in A_1\) such that, for every pair of integers \(i\geq0\) and \(r\geq1\),
\[
 \operatorname{rank}_k(\,\cdot\ell^r:A_i\longrightarrow A_{i+r})
 =\min\{\dim_k A_i,\dim_k A_{i+r}\}.
\]
After specifying the field, the sphere class, the neighborliness convention, and a property \(\mathcal P\), the proposed universal assertion takes the form
\[
 \text{for every admissible neighborly sphere }\Delta
 \text{ and every linear system of parameters }\Theta,
 \quad A(\Delta,\Theta)\text{ has }\mathcal P.
\]
In either displayed Lefschetz interpretation, the Lefschetz element may depend on \(\Delta\) and \(\Theta\). The source explicitly quantifies over every linear system of parameters, rather than only generic choices. It does not identify \(\mathcal P\) as weak or strong Lefschetz, specify characteristic restrictions, or determine whether the field is fixed or universally quantified. Consequently these candidate formulations must remain unresolved alternatives pending source clarification.
\subsection{Short English statement}
The question asks whether every choice of linear parameters for a neighborly sphere yields an Artinian quotient with the intended Lefschetz property. A precise statement still requires the coefficient field, the meaning of sphere and neighborliness, and whether weak or strong Lefschetz is intended. The supplied record labels the question open; this draft does not independently verify that status.
\subsection{Sources}
\begin{itemize}
\item Local UnsolvedMath registry, supplied complete record with permanent ID \(30004528\), source problem code \texttt{catalogue entry}; registry collection \texttt{data/unsolvedmath/problems.json}, with additional existing records maintained in \texttt{catalog/problems.json}. The supplied record identifies the set as Oberwolfach Reports Open Problems. No individual proposer is supplied.
\item Supplied source citation: \emph{Lefschetz Properties in Algebra, Geometry and Combinatorics} (2021), Oberwolfach Reports, pp.~1539--1600, DOI \texttt{10.4171/owr/2020/31}. Supplied source and literature link: \url{https://doi.org/10.4171/owr/2020/31}. The background describes these reports as documenting mathematical workshops and their open-problem discussions.
\item Supplied literature source: \emph{Lefschetz properties of spheres}, \url{https://arxiv.org/abs/0802.1058}. No authors or specific supporting passages are supplied in the record.
\item Literature assessment dated exactly \texttt{2026-08-22}: the supplied curation labels the question open and states that strong Lefschetz results cover important classes but do not establish the universal assertion. This is supplied dated curation, described in the record as literature triage derived from the OWR curation; neither the status nor the cited claims have been independently verified here.
\item The registry gives proposed year \(2020\), while its source citation labels the report \(2021\). Both supplied dates are retained without resolving their bibliographic relationship.
\end{itemize}
\end{document}
